% Options for packages loaded elsewhere
\PassOptionsToPackage{unicode}{hyperref}
\PassOptionsToPackage{hyphens}{url}
\documentclass[
]{article}
\usepackage{xcolor}
\usepackage{amsmath,amssymb}
\setcounter{secnumdepth}{-\maxdimen} % remove section numbering
\usepackage{iftex}
\ifPDFTeX
  \usepackage[T1]{fontenc}
  \usepackage[utf8]{inputenc}
  \usepackage{textcomp} % provide euro and other symbols
\else % if luatex or xetex
  \usepackage{unicode-math} % this also loads fontspec
  \defaultfontfeatures{Scale=MatchLowercase}
  \defaultfontfeatures[\rmfamily]{Ligatures=TeX,Scale=1}
\fi
\usepackage{lmodern}
\ifPDFTeX\else
  % xetex/luatex font selection
\fi
% Use upquote if available, for straight quotes in verbatim environments
\IfFileExists{upquote.sty}{\usepackage{upquote}}{}
\IfFileExists{microtype.sty}{% use microtype if available
  \usepackage[]{microtype}
  \UseMicrotypeSet[protrusion]{basicmath} % disable protrusion for tt fonts
}{}
\makeatletter
\@ifundefined{KOMAClassName}{% if non-KOMA class
  \IfFileExists{parskip.sty}{%
    \usepackage{parskip}
  }{% else
    \setlength{\parindent}{0pt}
    \setlength{\parskip}{6pt plus 2pt minus 1pt}}
}{% if KOMA class
  \KOMAoptions{parskip=half}}
\makeatother
\usepackage{longtable,booktabs,array}
\newcounter{none} % for unnumbered tables
\usepackage{calc} % for calculating minipage widths
% Correct order of tables after \paragraph or \subparagraph
\usepackage{etoolbox}
\makeatletter
\patchcmd\longtable{\par}{\if@noskipsec\mbox{}\fi\par}{}{}
\makeatother
% Allow footnotes in longtable head/foot
\IfFileExists{footnotehyper.sty}{\usepackage{footnotehyper}}{\usepackage{footnote}}
\makesavenoteenv{longtable}
\setlength{\emergencystretch}{3em} % prevent overfull lines
\providecommand{\tightlist}{%
  \setlength{\itemsep}{0pt}\setlength{\parskip}{0pt}}
% ===== unicode + compile header for ZIMMERMAN_LAW_OF_GRAVITY =====
% TO COMPILE ON OVERLEAF: upload this .tex + a figures/ folder containing
%   fig1_rar.png  fig2_a0z.png  fig3_btfr.png
% Compile with XeLaTeX (recommended) or pdfLaTeX. Download the PDF and publish
% manually on Zenodo. (Not test-compiled locally; unicode is mapped below.)
\usepackage{amssymb}
\usepackage{newunicodechar}
\providecommand{\textsubscript}[1]{\ensuremath{_{\mbox{\scriptsize #1}}}}
\providecommand{\textonehalf}{\ensuremath{\frac12}}
\newunicodechar{₀}{\ensuremath{_{0}}}
\newunicodechar{₁}{\ensuremath{_{1}}}
\newunicodechar{₃}{\ensuremath{_{3}}}
\newunicodechar{₆}{\ensuremath{_{6}}}
\newunicodechar{ₐ}{\ensuremath{_{a}}}
\newunicodechar{¹}{\ensuremath{^{1}}}
\newunicodechar{²}{\ensuremath{^{2}}}
\newunicodechar{⁰}{\ensuremath{^{0}}}
\newunicodechar{⁴}{\ensuremath{^{4}}}
\newunicodechar{⁵}{\ensuremath{^{5}}}
\newunicodechar{⁻}{\ensuremath{^{-}}}
\newunicodechar{Λ}{\ensuremath{\Lambda}}
\newunicodechar{ρ}{\ensuremath{\rho}}
\newunicodechar{π}{\ensuremath{\pi}}
\newunicodechar{σ}{\ensuremath{\sigma}}
\newunicodechar{Υ}{\ensuremath{\Upsilon}}
\newunicodechar{χ}{\ensuremath{\chi}}
\newunicodechar{Ω}{\ensuremath{\Omega}}
\newunicodechar{Δ}{\ensuremath{\Delta}}
\newunicodechar{γ}{\ensuremath{\gamma}}
\newunicodechar{α}{\ensuremath{\alpha}}
\newunicodechar{µ}{\ensuremath{\mu}}
\newunicodechar{√}{\ensuremath{\surd}}
\newunicodechar{×}{\ensuremath{\times}}
\newunicodechar{−}{\ensuremath{-}}
\newunicodechar{·}{\ensuremath{\cdot}}
\newunicodechar{∝}{\ensuremath{\propto}}
\newunicodechar{≈}{\ensuremath{\approx}}
\newunicodechar{≳}{\ensuremath{\gtrsim}}
\newunicodechar{≥}{\ensuremath{\geq}}
\newunicodechar{⇒}{\ensuremath{\Rightarrow}}
\newunicodechar{→}{\ensuremath{\rightarrow}}
\newunicodechar{↔}{\ensuremath{\leftrightarrow}}
\newunicodechar{⟺}{\ensuremath{\Longleftrightarrow}}
\newunicodechar{⋆}{\ensuremath{\star}}
\newunicodechar{★}{\ensuremath{\star}}
\newunicodechar{✓}{\ensuremath{\checkmark}}
\newunicodechar{✗}{\ensuremath{\times}}
\newunicodechar{◐}{\ensuremath{\circ}}
\newunicodechar{½}{\ensuremath{\frac{1}{2}}}
\newunicodechar{§}{\S{}}
\newunicodechar{Ö}{\"{O}}
\newunicodechar{Ü}{\"{U}}
\newunicodechar{é}{\'{e}}
\newunicodechar{³}{\ensuremath{^{3}}}
\newunicodechar{ä}{\"{a}}
\newunicodechar{Ȳ}{\ensuremath{\bar{Y}}}
\newunicodechar{δ}{\ensuremath{\delta}}
\newunicodechar{θ}{\ensuremath{\theta}}
\newunicodechar{κ}{\ensuremath{\kappa}}
\newunicodechar{λ}{\ensuremath{\lambda}}
\newunicodechar{μ}{\ensuremath{\mu}}
\newunicodechar{φ}{\ensuremath{\varphi}}
\newunicodechar{–}{\textendash{}}
\newunicodechar{—}{\textemdash{}}
\newunicodechar{…}{\ldots{}}
\newunicodechar{′}{'}
\newunicodechar{⁶}{\ensuremath{^{6}}}
\newunicodechar{⁷}{\ensuremath{^{7}}}
\newunicodechar{₂}{\ensuremath{_{2}}}
\newunicodechar{₄}{\ensuremath{_{4}}}
\newunicodechar{ℏ}{\ensuremath{\hbar}}
\newunicodechar{∇}{\ensuremath{\nabla}}
\newunicodechar{≔}{\ensuremath{:=}}
\newunicodechar{≤}{\ensuremath{\leq}}
\newunicodechar{≲}{\ensuremath{\lesssim}}
\newunicodechar{⊙}{\ensuremath{\odot}}
\newunicodechar{Σ}{\ensuremath{\Sigma}}
\newunicodechar{η}{\ensuremath{\eta}}
\newunicodechar{ν}{\ensuremath{\nu}}
\newunicodechar{″}{''}
\newunicodechar{⇔}{\ensuremath{\Leftrightarrow}}
\newunicodechar{±}{\ensuremath{\pm}}
\newunicodechar{¼}{\ensuremath{\tfrac{1}{4}}}
\newunicodechar{β}{\ensuremath{\beta}}
\newunicodechar{∈}{\ensuremath{\in}}
\newunicodechar{∞}{\ensuremath{\infty}}
\newunicodechar{≠}{\ensuremath{\neq}}
% --- v3 erratum build: additional symbol coverage ---
\newunicodechar{ω}{\ensuremath{\omega}}
\newunicodechar{τ}{\ensuremath{\tau}}
\newunicodechar{∫}{\ensuremath{\int}}
\newunicodechar{∂}{\ensuremath{\partial}}
\newunicodechar{≫}{\ensuremath{\gg}}
\newunicodechar{⟨}{\ensuremath{\langle}}
\newunicodechar{⟩}{\ensuremath{\rangle}}
\newunicodechar{⁸}{\ensuremath{^{8}}}
\usepackage{bookmark}
\IfFileExists{xurl.sty}{\usepackage{xurl}}{} % add URL line breaks if available
\urlstyle{same}
\hypersetup{
  hidelinks,
  pdfcreator={LaTeX via pandoc}}

\author{}
\date{}

\begin{document}

\section{\texorpdfstring{A Pre-Registered \(a_0(z)\) Gate for the
Rubin/LSST Supernova Stream: Committing the Test Before the
Data}{A Pre-Registered a\_0(z) Gate for the Rubin/LSST Supernova Stream: Committing the Test Before the Data}}\label{a-pre-registered-a_0z-gate-for-the-rubinlsst-supernova-stream-committing-the-test-before-the-data}

\textbf{Carl P. Zimmerman} Briar Creek Tech --- carl@briarcreektech.com
Frozen 2026-07-21

\begin{center}\rule{0.5\linewidth}{0.5pt}\end{center}

\subsection{Abstract}\label{abstract}

\textbf{This is a pre-registration, not a result.} It commits --- in
advance, timestamped, and SHA-256-hashed --- the complete analysis
procedure by which one distinctive prediction of the de Sitter--Unruh
\emph{modified-inertia} framework will be judged against the forthcoming
Rubin Observatory / LSST supernova cosmology stream. The framework ties
the galaxy acceleration scale \(a_0\) to the cosmological horizon and,
under an adiabatic-horizon ansatz, predicts that \(a_0\) tracks the
square root of the dark-energy density:
\(a_0(z)/a_0(0)=\sqrt{\rho_\mathrm{DE}(z)/\rho_\mathrm{DE,0}}\). The
observable is the ratio \(R\equiv a_0(z{=}3)/a_0(0)\), which is
\emph{independent of the normalization} \(Z\) and of the absolute value
of \(a_0\) --- both footings of the framework collapse onto the same
\(R\). The distinctive claim is that \(R\) declines to
\(\approx 0.60\)--\(0.75\) \textbf{if and only if} dark energy evolves;
if \(w=-1\), then \(R=1\) and the framework degenerates to
constant-\(a_0\) MOND. We freeze the estimator, the decision thresholds,
and the data hierarchy \emph{before} any calibrated Rubin supernova
cosmology exists --- LSST began operations on 2026-06-30, and the
LSST-DESC calibrated Type-Ia sample is not expected before
\textasciitilde2027. The credibility move is precisely that the test is
written before the data can shape it. At the freeze the best available
input (DESI DR2 2025) gives \(R=0.775\,[0.68,0.88]\), a \(2.0\sigma\)
decline --- below the committed \(3\sigma\) bar, hence
\textbf{UNDECIDED}. A committed forecast projects that Rubin sharpens
this to \(\sim 3.3\sigma\) (supernovae alone) up to \(\sim 4.7\sigma\)
(supernovae combined with CMB/BAO) if dark energy evolves at the
DESI-central rate. We state explicitly the \emph{unsound path we
refuse}: no do-it-yourself Hubble diagram from public broker light
curves. Every number here comes from a committed, hash-locked script;
the Zenodo DOI and timestamp are the public lock-in.

\begin{center}\rule{0.5\linewidth}{0.5pt}\end{center}

\subsection{1. Introduction}\label{introduction}

\subsubsection{1.1 The prediction under
test}\label{the-prediction-under-test}

The de Sitter--Unruh modified-inertia framework is built on a
horizon-derived acceleration scale \(a_0 = cH_\Lambda/Z\), with the
canonical footing giving \(a_0 = 9.36\times10^{-11}\,{\rm m\,s^{-2}}\)
and an alternate footing giving \(1.13\times10^{-10}\,{\rm m\,s^{-2}}\).
The low-acceleration kernel it uses, \(\nu = \sqrt{1+1/y}\), is the
identical functional form to Milgrom's 1999 interpolation (Milgrom 1983,
1999); the framework's distinctive content is the \emph{coefficient} ---
the tie of the scale to \(cH_\Lambda/Z\) rather than to a free constant
--- together with a modified-inertia (rather than modified-gravity)
completion.

If the acceleration scale is sourced by the cosmological horizon, then
when the horizon evolves, so should \(a_0\). Under the adiabatic-horizon
ansatz the scale tracks the dark-energy density: \[
a_0(z)/a_0(0) = \sqrt{\rho_\mathrm{DE}(z)/\rho_\mathrm{DE,0}}, \qquad
\frac{\rho_\mathrm{DE}(z)}{\rho_\mathrm{DE,0}} = (1+z)^{3(1+w_0+w_a)}\, e^{-3 w_a z/(1+z)},
\] where the second equality is the standard CPL
(Chevallier--Polarski--Linder) parameterization of an evolving equation
of state \(w(a)=w_0+w_a(1-a)\). The single scalar observable is \[
R \equiv a_0(z{=}3)/a_0(0).
\] That \(a_0\) should scale as \(\sqrt{\rho_\mathrm{DE}}\) is not original
to this framework --- Limbach, Psaltis \& Özel (2008) noted the
numerical coincidence \(a_0\sim c\sqrt{\rho_\mathrm{DE}}\) well before it.
What the framework adds is the \emph{promotion} of that coincidence to a
horizon-sourced dynamical law, which is why it makes a falsifiable
redshift prediction rather than a static numerical match.

\textbf{\(R\) is \(Z\)-independent.} Because \(R\) is a \emph{ratio} of
\(a_0\) at two redshifts, the normalization \(Z\) and the absolute value
of \(a_0\) cancel identically. Both \emph{value} footings --- canonical
\(9.36\times10^{-11}\) and alternate \(1.13\times10^{-10}\) --- collapse
onto the same \(R(w_0,w_a)\). This is what makes the gate a clean
cosmology-side test: it depends on the dark-energy expansion history
alone, not on any posited normalization internal to the framework. One
clarification is owed here, because the framework's internal audit
distinguishes two separate footing forks. The one that cancels is the
\emph{value} fork (which absolute \(a_0\) today). A \emph{second},
logically independent fork concerns the \emph{evolution law itself} ---
the adiabatic-horizon \(a_0\propto\sqrt{\rho_\mathrm{DE}}\) (which declines
under evolving DE) versus a rival \(a_0\propto cH(z)=cH_0E(z)\) (which
\emph{rises} with \(z\)). \textbf{This pre-registration commits to the
\(\sqrt{\rho_\mathrm{DE}}\) footing as the prediction under test.} The
rival rising-\(cH E(z)\) footing is not adopted here; if the data
instead showed a rise, that registers as Verdict C below. So ``both
footings collapse to \(R\)'' is a statement about the value
normalization only --- it is not a claim that the evolution law is
unique.

The distinctive claim is:

\begin{quote}
\(R\) declines to \(\approx 0.60\)--\(0.75\) \textbf{if and only if}
dark energy evolves. If \(w=-1\) (a true cosmological constant), then
\(R=1\) exactly and the framework reduces to ordinary constant-\(a_0\)
MOND.
\end{quote}

\textbf{Two facts about the shape of \(a_0(z)\) must be stated up front,
because the gate summarizes a whole curve by a single point.} First,
under an evolving CPL history \(a_0(z)\) is \emph{not monotonic}: at the
DESI-central rate it rises to a small \textbf{bump} --- \(R\) peaks at
\(\approx1.036\) near \(z\approx0.35\) and only crosses back below unity
at \(z\approx0.9\) --- and \emph{declines} only at higher redshift. This
is the framework's own committed ``bump-then-decline'' behavior, which
the estimator's \(\sqrt{\rho_\mathrm{DE}}\) form reproduces identically.
Second, the gate is evaluated at a \textbf{frozen reference redshift
\(z=3\)}, deliberately beyond the range the feeding supernovae actually
cover (\(z\lesssim1.2\) for the LSST photometric sample).
\(R=a_0(3)/a_0(0)\) is therefore a \emph{CPL extrapolation}: the
headline decline is a lever-arm statement, not a directly measured
\(z=3\) quantity. Operationally, the significance the gate reports is
the significance with which the supernova-constrained \((w_0,w_a)\)
posterior excludes the non-evolving point \((w_0,w_a)=(-1,0)\), read
along the CPL curve at \(z=3\). Both facts are quantified and their
consequences drawn out in §6.1.

\subsubsection{1.2 Why pre-register now, before Rubin data
exist}\label{why-pre-register-now-before-rubin-data-exist}

The Vera C. Rubin Observatory began its Legacy Survey of Space and Time
on 2026-06-30. No calibrated Rubin supernova cosmology sample yet
exists; the LSST-DESC systematics-controlled Type-Ia sample is not
expected before roughly 2027. This is exactly the window in which a
pre-registration has value. Committing the estimator and the decision
thresholds \emph{before} the data land removes any later freedom to tune
the analysis to whatever the data happen to show --- the single largest
source of after-the-fact ``detections'' in cosmology-adjacent model
testing.

This follows the same discipline as the author's Gaia-DR4
pre-registration precedent (frozen wide-binary \(\gamma\) targets and
cut tables, filed before DR4), now applied to Rubin. The move is
deliberately modest: we are not announcing a finding, we are locking the
referee-grade procedure so that when the data arrive there is nothing
left to negotiate but the input \((w_0,w_a)\) posterior.

\subsubsection{1.3 What this pre-registration does and does not
settle}\label{what-this-pre-registration-does-and-does-not-settle}

This gate settles the \textbf{cosmology-side} question --- \emph{does
\(a_0\) evolve at all} --- which the framework \emph{inherits} from the
dark-energy history \(w(z)\). It is important to state at the outset,
and it is frozen into the protocol below, that a favorable verdict means
the distinctive prediction \textbf{survives}, not that the framework is
established. The independent, corroborating test is the
\textbf{galaxy-side} measurement of \(a_0(z)\) from resolved
rotation-curve dynamics across redshift (the cross-scale program, DOI
10.5281/zenodo.21440407). The value of \(a_0\), the normalization \(Z\),
and the adiabatic-horizon promotion of
\(a_0(z)\propto\sqrt{\rho_\mathrm{DE}}\) are all \textbf{posited}, not
derived. This scope statement is non-negotiable and appears again,
verbatim, in the frozen estimator's docstring.

\begin{center}\rule{0.5\linewidth}{0.5pt}\end{center}

\subsection{2. The Frozen Estimator and Pre-Committed
Thresholds}\label{the-frozen-estimator-and-pre-committed-thresholds}

The decision procedure is committed in \texttt{a0z\_gate\_estimator.py}
(SHA-256 in §8). \textbf{Input:} a calibrated \((w_0,w_a)\) posterior
--- mean, marginal uncertainties, and correlation coefficient.
\textbf{Method:} the estimator draws \(N=4\times10^5\) samples from the
input's bivariate Gaussian (via Cholesky factorization of the
covariance), maps each draw through the framework relation to
\(R=a_0(3)/a_0(0)\), and reports the median with its \([16,84]\)
percentile interval, the significance of a \emph{decline} (\(R<1\))
against flat, and the significance of a \emph{rise} (\(R>1\)). It then
applies the following pre-committed thresholds. \textbf{No threshold,
band, or estimator may change after the freeze; only the input
\((w_0,w_a)\) is updated.}

{\def\LTcaptype{none} % do not increment counter
\begin{longtable}[]{@{}
  >{\raggedright\arraybackslash}p{(\linewidth - 2\tabcolsep) * \real{0.5000}}
  >{\raggedright\arraybackslash}p{(\linewidth - 2\tabcolsep) * \real{0.5000}}@{}}
\toprule\noalign{}
\begin{minipage}[b]{\linewidth}\raggedright
Verdict
\end{minipage} & \begin{minipage}[b]{\linewidth}\raggedright
Pre-committed condition
\end{minipage} \\
\midrule\noalign{}
\endhead
\bottomrule\noalign{}
\endlastfoot
\textbf{A) GATE OPEN} --- distinctive prediction alive and supported &
decline detected at \(\ge 3\sigma\) \textbf{and} \(R\in[0.55,0.85]\)
(the \(\sqrt{\rho_\mathrm{DE}}\) band across both footings) → \(a_0\)
evolves as required; warrants the independent galaxy-side
confrontation \\
\textbf{B) GATE DISSOLVED} --- safe core, \emph{not} a falsification &
consistent with flat (\(R=1\) within \(2\sigma\)) → framework reduces to
constant-\(a_0\) MOND; the distinctive prediction is inapplicable. If in
addition \(R\le0.85\) is excluded at \(\ge3\sigma\), the distinctive
decline is positively ruled out (still dissolution, not a kill of the
MOND-scale core) \\
\textbf{C) FRAMEWORK STRAINED} --- cosmology-side tension & \(R>1\) (a
\emph{rise}) at \(\ge3\sigma\) → the scale rises with \(z\), which
\(a_0\propto\sqrt{\rho_\mathrm{DE}}\) does not permit under evolving DE \\
\textbf{UNDECIDED} & none of the above met (insufficient significance /
out of band) \\
\end{longtable}
}

The band \([0.55,0.85]\) for Verdict A is the range of
\(R=a_0(3)/a_0(0)\) produced by the \(\sqrt{\rho_\mathrm{DE}}\) relation
across the range of dark-energy histories currently plausible; it is
fixed here and does not move. To be precise about what sets it --- and
to forestall any charge of gerrymandering --- the band width comes from
the spread of allowed \((w_0,w_a)\), \textbf{not} from the two value
footings (those give an \emph{identical} \(R\) at fixed \((w_0,w_a)\),
so they cannot widen a band). Holding \(w_a=-0.62\), the edges
correspond to \(w_0\approx-0.79\) (\(R=0.85\), mild evolution) and
\(w_0\approx-1.00\) (\(R=0.55\), strong evolution). Two honesty notes
follow. (i) The band \([0.55,0.85]\) is deliberately \emph{wider} than
the nominal distinctive claim \(R\approx0.60\)--\(0.75\) quoted in §1.1;
the DESI-central value \(R=0.775\) in fact sits just \emph{above} the
nominal \(0.75\) but inside the operational band. (ii) The upper edge
\(R=0.85\) is close to flat, so a ``supported'' verdict at that edge is
a weak decline; this is the intended cost of choosing a decision band
broad enough not to be tuned to any single posterior. The asymmetry
between Verdicts A, B, and C is intentional: a decline in-band supports
the distinctive claim, flatness dissolves it into the safe MOND core,
and a \emph{rise} would strain the framework because the committed
\(\sqrt{\rho_\mathrm{DE}}\) footing produces no rising \(a_0\) under
evolving dark energy. Note the deliberate gap: a decline so steep that
\(R<0.55\) at \(\ge3\sigma\) (steeper than \(\sqrt{\rho_\mathrm{DE}}\)
allows) returns \textbf{UNDECIDED}, not a verdict --- the frozen
procedure has no ``over-decline strain'' branch mirroring Verdict C (see
§6.1).

\begin{center}\rule{0.5\linewidth}{0.5pt}\end{center}

\subsection{3. Data Hierarchy and the Refused DIY-Hubble-Diagram
Path}\label{data-hierarchy-and-the-refused-diy-hubble-diagram-path}

The gate is fed by a strict, frozen hierarchy of inputs:

\begin{enumerate}
\def\labelenumi{\arabic{enumi}.}
\tightlist
\item
  \textbf{Primary:} the LSST-DESC \emph{calibrated} Type-Ia supernova
  cosmology sample --- systematics-controlled zero-points,
  photometric-redshift treatment, and a modeled selection function.
  Expected \textasciitilde2027 and later. This is the only input that
  can produce a \emph{supported} verdict.
\item
  \textbf{Interim:} published Rubin \(w_0w_a\) constraints as they
  appear, including externally combined CMB/BAO fits. These update the
  input posterior but carry the caveats of any early, pre-final
  calibration.
\item
  \textbf{Readiness tracker only:} a running count of broker-classified
  Type-Ia candidates from the \emph{public} alert stream
  (\texttt{broker\_sn\_counter.py}). \textbf{This count never feeds the
  verdict.} It tracks \emph{when} the gate will have enough statistical
  weight to spring --- timing, not cosmology.
\end{enumerate}

\subsubsection{The path we refuse}\label{the-path-we-refuse}

We do \textbf{not} build a do-it-yourself Hubble diagram from public
broker light curves. The alerts are public and real-time, but they are
discoveries, not calibrated distances. Without the DESC calibration
layer --- zero-point stability, photometric-redshift systematics, and a
modeled selection function --- any \(w(z)\) extracted from raw broker
photometry is systematics-dominated, and a ``Rubin supports \(a_0(z)\)''
claim built on it would be a manufactured result. Refusing that path is
stated here as a virtue of the protocol, not an omission: the primary
input is the calibrated sample, full stop. The readiness tracker exists
precisely so that the temptation to jump early has a disciplined,
cosmology-free outlet.

\begin{center}\rule{0.5\linewidth}{0.5pt}\end{center}

\subsection{4. Forecast: What Significance Rubin
Reaches}\label{forecast-what-significance-rubin-reaches}

The forecast basis is committed in \texttt{forecast\_rubin\_a0z.py}
(SHA-256 in §8); it is explicitly a projection, not a measurement. It
propagates representative \emph{projected} Rubin \(w_0\)--\(w_a\)
covariances through the framework relation and asks at what significance
Rubin distinguishes the framework's decline from flat (\(w=-1\)).
Photometric supernova cosmology forecasts for LSST reach a Dark Energy
Task Force figure of merit of order 150 (LSST alone) up to
\textasciitilde500 (LSST combined with CMB/BAO), against
\textasciitilde54 for the Dark Energy Survey and the
\textasciitilde{}\(2.0\sigma\) decline currently delivered by DESI DR2.
Taking the DESI-central evolving-DE history as the assumed truth
(\(w_0=-0.838\), \(w_a=-0.62\)), the committed forecast gives:

{\def\LTcaptype{none} % do not increment counter
\begin{longtable}[]{@{}
  >{\raggedright\arraybackslash}p{(\linewidth - 8\tabcolsep) * \real{0.2000}}
  >{\raggedright\arraybackslash}p{(\linewidth - 8\tabcolsep) * \real{0.2000}}
  >{\raggedright\arraybackslash}p{(\linewidth - 8\tabcolsep) * \real{0.2000}}
  >{\raggedright\arraybackslash}p{(\linewidth - 8\tabcolsep) * \real{0.2000}}
  >{\raggedright\arraybackslash}p{(\linewidth - 8\tabcolsep) * \real{0.2000}}@{}}
\toprule\noalign{}
\begin{minipage}[b]{\linewidth}\raggedright
Input scenario
\end{minipage} & \begin{minipage}[b]{\linewidth}\raggedright
\(\sigma(w_0)\)
\end{minipage} & \begin{minipage}[b]{\linewidth}\raggedright
\(\sigma(w_a)\)
\end{minipage} & \begin{minipage}[b]{\linewidth}\raggedright
\(R=a_0(3)/a_0(0)\)
\end{minipage} & \begin{minipage}[b]{\linewidth}\raggedright
Decline significance
\end{minipage} \\
\midrule\noalign{}
\endhead
\bottomrule\noalign{}
\endlastfoot
DESI DR2 now (baseline) & 0.055 & 0.22 & 0.775 {[}0.683, 0.878{]} &
\(2.0\sigma\) \\
Rubin SN alone (FoM \textasciitilde150) & 0.040 & 0.15 & 0.775 {[}0.717,
0.837{]} & \(3.3\sigma\) \\
Rubin SN + CMB/BAO (FoM \textasciitilde500) & 0.020 & 0.08 & 0.775
{[}0.742, 0.809{]} & \(4.6\)--\(4.7\sigma\) \\
\end{longtable}
}

The forecast also runs the \textbf{mirror test} --- if the universe is
in fact flat (\(w=-1\), \(a_0\) truly constant), how tightly does Rubin
pin \(R\) to unity and thereby exclude the framework's \(\approx0.74\)?
Rubin-alone (FoM \textasciitilde150) gives \(R=1.00\pm0.078\), excluding
\(0.74\) at \(3.3\sigma\); Rubin + CMB/BAO (FoM \textasciitilde500)
gives \(R=1.00\pm0.043\), excluding it at \(6.0\sigma\).

The honest forecast summary is deliberately unheroic: if dark energy
evolves at the DESI-central rate, Rubin sharpens the \(a_0\)-decline
detection from DESI-now \(\sim2.0\sigma\) to \(\sim3.3\sigma\)
(supernovae alone) or \(\sim4.6\)--\(4.7\sigma\) (with CMB/BAO) ---
crossing the committed \(3\sigma\) bar, approaching but not alone
reaching \(5\sigma\) at the modest \(0.775\)-vs-\(1.0\) gap. If dark
energy is flat, Rubin pins \(R=1\) and excludes the framework's \(0.74\)
at \(\sim3.3\)--\(6.0\sigma\). Either way the gate is likely settled at
the \(3\)--\(6\sigma\) level by roughly 2028--2030 --- a real upgrade on
DESI-now, not a one-instrument slam dunk.

\begin{center}\rule{0.5\linewidth}{0.5pt}\end{center}

\subsection{5. State at the Freeze
(2026-07-21)}\label{state-at-the-freeze-2026-07-21}

Running the frozen estimator on the best current input reproduces the
baseline verdict exactly:

\begin{itemize}
\item
  \textbf{Cosmology input (DESI DR2 2025, arXiv:2503.14738):}
  \(w_0=-0.838\pm0.055\), \(w_a=-0.62\pm0.22\), correlation \(-0.86\) →
  \(R=0.775\,[0.683,0.878]\), a decline at \(2.0\sigma\). This is
  \textbf{below the committed \(3\sigma\) threshold}, so the frozen
  verdict at the freeze is \textbf{UNDECIDED} --- a \(2.0\sigma\) hint,
  explicitly not a detection, explicitly not gate-open. Applying the
  same procedure to the two forecast inputs returns \emph{GATE OPEN} at
  \(3.3\sigma\) and \(4.7\sigma\) respectively, showing that the
  projected Rubin precision is what moves the verdict off UNDECIDED ---
  but those are projections, not data, and cannot produce a real
  verdict.
\item
  \textbf{Broker readiness tracker (live, 2026-07-21):} the ALeRCE
  light-curve classifier reports \textbf{36,836} Type-Ia-classified
  objects (ZTF legacy plus early LSST). This is past the
  \textasciitilde4,400 spectroscopic-quality threshold
  (\textbf{REACHED}) and is \textbf{9.21\%} of the
  \textasciitilde400,000 LSST photometric target. It is a
  broker-classifier count --- unvetted, no redshifts, ZTF-and-early-LSST
  mixed --- and by the frozen hard rule it \textbf{never feeds the
  verdict}. It reports timing only.
\end{itemize}

The fallback brokers, should the ALeRCE endpoint be unavailable, are
Fink, Lasair, ANTARES, and Pitt-Google; the tracker script is
ready-to-run against any of them.

\begin{center}\rule{0.5\linewidth}{0.5pt}\end{center}

\subsection{6. Honest Scope, and the Kill / Dissolution
Conditions}\label{honest-scope-and-the-kill-dissolution-conditions}

The scope is frozen in and stated plainly:

\begin{itemize}
\tightlist
\item
  This settles the \textbf{cosmology-side gate} --- \emph{does \(a_0\)
  evolve at all} --- which the framework inherits from \(w(z)\).
  \textbf{GATE OPEN means the distinctive prediction survives, not that
  the framework is established.} The corroborating test remains the
  \emph{independent galaxy-side} \(a_0(z)\) measured from resolved
  dynamics across redshift (cross-scale, DOI 10.5281/zenodo.21440407).
  The cosmology-side gate and the galaxy-side test are logically
  distinct; passing the former is necessary, not sufficient.
\item
  \(a_0\)'s value, the normalization \(Z\), and the adiabatic-horizon
  promotion of \(a_0(z)\propto\sqrt{\rho_\mathrm{DE}}\) are
  \textbf{posited}, not derived. The Milgrom-form kernel
  \(\nu=\sqrt{1+1/y}\) is credited to Milgrom (1999); the framework's
  distinctive content is the horizon coefficient and the
  modified-inertia completion (MI Field Theory, DOI
  10.5281/zenodo.21403470).
\end{itemize}

The pre-committed kill / dissolution conditions, signed 2026-07-21:

\begin{itemize}
\tightlist
\item
  \textbf{Distinctive prediction FALSIFIED} if the calibrated Rubin
  posterior yields \textbf{Verdict C} --- \(a_0\) \emph{rises} with
  \(z\) at \(\ge3\sigma\) --- which no viable footing of
  \(a_0\propto\sqrt{\rho_\mathrm{DE}}\) permits under evolving dark energy.
\item
  \textbf{Distinctive prediction DISSOLVED} (Verdict B) if \(w\to-1\):
  the framework becomes ordinary constant-\(a_0\) MOND. This is the
  stated safe core --- dissolution, not a kill of the MOND-scale
  reframing. The reframing survives as constant-\(a_0\) MOND even if the
  redshift signature vanishes.
\item
  \textbf{Distinctive prediction SUPPORTED} (Verdict A) only if
  \emph{both} the \(\ge3\sigma\) decline \emph{and} the \([0.55,0.85]\)
  band are met --- and even then, ``supported'' means survives to face
  the galaxy-side test, nothing more.
\end{itemize}

Because \(R\) is \(Z\)-independent, none of these outcomes depends on
the internal normalization; they depend on the measured dark-energy
history alone. This is the cleanest form the test can take, and also the
sharpest constraint on what a favorable outcome is allowed to mean.

\subsubsection{6.1 Known limitations baked into the frozen
procedure}\label{known-limitations-baked-into-the-frozen-procedure}

Because the estimator and thresholds are hashed and cannot be revised
after the freeze, the following limitations are documented rather than
fixed. They are stated here so that no future reading can present the
gate as tighter than it is.

\begin{enumerate}
\def\labelenumi{\arabic{enumi}.}
\item
  \textbf{The signal is a \(z=3\) extrapolation beyond the supernova
  range, and the significance is \(z\)-dependent.} The LSST photometric
  SN sample constrains \(w(z)\) mainly at \(z\lesssim1.2\); the gate is
  read at \(z=3\). Propagating the \emph{same} DESI-DR2 posterior to
  different reference redshifts gives a decline significance of
  \(0.34\sigma\) at \(z=1\) (\(R=0.99\), effectively flat),
  \(1.7\sigma\) at \(z=2\), \(2.0\sigma\) at \(z=3\), and \(2.4\sigma\)
  at \(z=10\); for the tight SN+CMB/BAO forecast covariance it is
  \(1.1\sigma\) at \(z=1\) but \(4.7\sigma\) at \(z\ge2\). The headline
  numbers therefore depend on the frozen choice \(z=3\), which is a
  lever arm, not a redshift at which calibrated SNe exist. What the gate
  genuinely tests is whether the SN-constrained posterior excludes
  \((w_0,w_a)=(-1,0)\); the \(z=3\) evaluation converts that into an
  \(a_0\)-ratio but does not add independent \(z=3\) information. A
  reader must not treat \(R(z{=}3)=0.775\) as a measured \(z=3\)
  acceleration scale.
\item
  \textbf{The framework's own \(a_0(z)\) bumps \emph{up} at low
  redshift, which is where the galaxy-side corroborating test lives.} At
  the DESI-central history \(a_0(z)\) rises to \(R\approx1.036\) near
  \(z\approx0.35\) and only falls below \(1\) beyond \(z\approx0.9\).
  The independent galaxy-side \(a_0(z)\) program (resolved
  rotation-curve dynamics, DOI 10.5281/zenodo.21440407) is accessible
  only at low redshift --- squarely inside this \emph{bump}. So a
  GATE-OPEN verdict on the \(z=3\) \emph{decline} does not translate
  into a low-\(z\) decline for the galaxy-side to corroborate; at
  \(z\lesssim0.5\) the framework predicts a slight \textbf{rise} (a few
  percent), and at \(z\sim1\) essentially no change. Anyone carrying
  ``the framework predicts \(a_0\) declines'' to galaxy-side observers
  at accessible redshifts would have the sign wrong. The two halves of
  the program probe opposite branches of the same non-monotonic curve,
  and the amplitude available to the galaxy side (\(\sim\)few percent)
  is far smaller than the \(z=3\) headline.
\item
  \textbf{Verdict C (``rise ⇒ strained'') is well-posed only at the
  frozen \(z=3\).} Because \(a_0(z)\) genuinely rises for
  \(z\lesssim0.9\), a measured rise is a contradiction of the
  \(\sqrt{\rho_\mathrm{DE}}\) footing \emph{only} when read at high \(z\);
  the frozen \(z=3\) evaluation is what makes the rule self-consistent.
  This is a reason the reference redshift is fixed, but it also means
  Verdict C should never be applied to a low-\(z\) input.
\item
  \textbf{Doc-vs-code items.} (a) The Verdict-B sub-clause described in
  §2 and §6 --- ``if \(R\le0.85\) is additionally excluded at
  \(\ge3\sigma\), the decline is positively ruled out'' --- is
  \emph{narrative only}; the frozen \texttt{a0z\_gate\_estimator.py}
  computes Verdict B purely as consistency with \(R=1\) within
  \(2\sigma\) and does not evaluate that exclusion. (b) The over-decline
  gap of §2 (\(R<0.55\) significant \(\to\) UNDECIDED) is a real
  asymmetry with Verdict C. (c) The frozen estimator prints ``GATE
  OPEN'' for the two \emph{forecast} covariances; those are fabricated
  inputs shown for illustration and, by the data hierarchy of §3, can
  never constitute a verdict --- only a calibrated posterior can.
\item
  \textbf{Significance is a Bayesian posterior tail, and the two frozen
  scripts differ at the \(0.1\sigma\) level.} The reported
  ``\(N\sigma\)'' is the posterior credibility that \(R<1\) (the
  fraction of samples with \(R\ge1\) mapped to a Gaussian sigma), not a
  frequentist likelihood-ratio detection. The estimator
  (\(N=4\times10^5\)) and the forecast script (\(N=3\times10^5\)) return
  \(4.7\sigma\) and \(4.6\sigma\) respectively for the identical
  SN+CMB/BAO scenario --- a Monte-Carlo difference, reported here as the
  range \(4.6\)--\(4.7\sigma\).
\item
  \textbf{DESI input provenance.} The baseline
  \((w_0,w_a)=(-0.838,-0.62)\), \(\sigma=(0.055,0.22)\), \(\rho=-0.86\)
  is the DESI DR2 + CMB + Pantheon+ combination (the mildest of the DESI
  DR2 SN combinations; DESY5 and Union3 prefer stronger evolution and a
  lower \(R\)). The verdict at the freeze is UNDECIDED under this input
  and remains UNDECIDED under the stronger combinations, so the choice
  does not change the frozen state --- but the input row is named
  explicitly so the update path is unambiguous.
\end{enumerate}

\begin{center}\rule{0.5\linewidth}{0.5pt}\end{center}

\subsection{7. Conclusion}\label{conclusion}

We have frozen, before the data exist, a referee-grade decision
procedure for the one distinctive prediction the de Sitter--Unruh
modified-inertia framework makes on the cosmological side: that the
galaxy acceleration scale tracks \(\sqrt{\rho_\mathrm{DE}}\) and therefore
\emph{declines} with redshift if and only if dark energy evolves. The
observable \(R=a_0(3)/a_0(0)\) is normalization-independent; the
estimator, the thresholds (\(3\sigma\) decline and the \([0.55,0.85]\)
band for GATE OPEN), the data hierarchy, and the refusal of the
DIY-Hubble-diagram path are all committed and hashed. At the freeze the
test stands \textbf{UNDECIDED}: DESI DR2 gives a \(2.0\sigma\) hint,
\(R=0.775\,[0.68,0.88]\), below the committed bar. The committed
forecast projects Rubin to reach \(\sim3.3\)--\(4.7\sigma\) if dark
energy evolves, or to exclude the decline at \(\sim3.3\)--\(6.0\sigma\)
if it does not --- a decision at the \(3\)--\(6\sigma\) level expected
by roughly 2028--2030 from the LSST-DESC calibrated sample. A favorable
outcome would mean the distinctive prediction survives to face the
independent galaxy-side \(a_0(z)\) test; it would not, on its own,
establish the framework. The value of this document is entirely in its
timing: it is filed before the calibrated data can shape it, and only
the \((w_0,w_a)\) input will be updated when they arrive.

\begin{center}\rule{0.5\linewidth}{0.5pt}\end{center}

\subsection{8. Freeze Integrity}\label{freeze-integrity}

The SHA-256 hashes of the frozen artifacts, recorded in
\texttt{FREEZE\_HASHES.txt} and verified 2026-07-21:

\begin{verbatim}
852ca954431b9239415038dde9f2b2ad9f89fe6325ff7a8ba3e3ebe357366740  a0z_gate_estimator.py
ba24927ddb3263c390646019b1dba1968c4aa3de5ed572bc5de57685afd2b77f  broker_sn_counter.py
bd714c93edd3eadc6c84413dee331e86be7ebaca8f0a5f88c84ec1542f46df0f  PREREGISTRATION.md
530e66beef425bb500d085051c3219669e6d6e31364eec9263ea036bca7a5697  forecast_rubin_a0z.py
\end{verbatim}

The Zenodo DOI and its deposit timestamp are the public lock-in of this
pre-registration. No threshold, band, or estimator changes after this
hash; only the calibrated \((w_0,w_a)\) input is updated when the
LSST-DESC sample lands.

\begin{center}\rule{0.5\linewidth}{0.5pt}\end{center}

\subsection{References}\label{references}

\begin{itemize}
\tightlist
\item
  DESI Collaboration (2025). \emph{DESI DR2 Results II: Measurements of
  Baryon Acoustic Oscillations and Cosmological Constraints.}
  arXiv:2503.14738; Phys. Rev.~D 112, 083515. (The
  \(w_0=-0.838,\,w_a=-0.62\) input used here is the DESI+CMB+Pantheon+
  combination.)
\item
  Ivezić, Ž., et al.~(LSST/Rubin) (2019). \emph{LSST: From Science
  Drivers to Reference Design and Anticipated Data Products.} ApJ 873,
  111.
\item
  LSST Dark Energy Science Collaboration (LSST-DESC). Supernova
  cosmology forecasts (photometric FoM \textasciitilde150; +CMB/BAO
  \textasciitilde500).
\item
  Möller, A., et al.~\emph{Fink}; Smith, K. W., et al.~\emph{Lasair};
  Matheson, T., et al.~\emph{ANTARES}; Förster, F., et al.~\emph{ALeRCE}
  --- LSST/ZTF alert brokers.
\item
  Milgrom, M. (1983). \emph{A modification of the Newtonian dynamics as
  a possible alternative to the hidden mass hypothesis.} ApJ 270, 365.
\item
  Milgrom, M. (1999). \emph{The modified dynamics as a vacuum effect.}
  Phys. Lett. A 253, 273 --- the \(\nu=\sqrt{1+1/y}\) kernel.
\item
  Limbach, C., Psaltis, D., Özel, F. (2008). \emph{The Redshift
  Evolution of the Tully--Fisher Relation as a Test of Modified
  Gravity.} arXiv:0809.2790 --- considered \(a_0\) coupled to \(H_0\) or
  to \(\rho_\mathrm{DE}\) (the \(a_0\sim c\sqrt{\rho_\mathrm{DE}}\)
  coincidence predates this framework) and tested both couplings against
  Tully--Fisher data to \(z\simeq1.2\). Their credit is for the
  coincidence; the horizon-sourced dynamical promotion is what this
  framework adds.
\item
  Zimmerman, C. P. \emph{Cross-scale \(a_0(z)\) program} (galaxy-side
  corroborating test). DOI 10.5281/zenodo.21440407.
\item
  Zimmerman, C. P. \emph{Modified-Inertia Field Theory.} DOI
  10.5281/zenodo.21403470.
\item
  Zimmerman, C. P. \emph{Gaia-DR4 pre-registration} (frozen-cuts
  precedent).
\end{itemize}

\begin{center}\rule{0.5\linewidth}{0.5pt}\end{center}

\emph{Author: Carl P. Zimmerman, Briar Creek Tech
(carl@briarcreektech.com). This is a pre-registration filed before the
calibrated Rubin supernova cosmology sample exists. It is a committed
analysis protocol, not a result. Every quantitative figure herein is
produced by a committed, hash-locked script that exits 0.}

\end{document}
